\chapter{Finite Temperature Field Theory: General formalism}
\label{ch:finite-temperature}

Conventional quantum field theory is formulated in empty space. The LSZ
formalism expresses scattering matrix elements in terms of zero-temperature
Green's functions, that is, vacuum expectation values of time-ordered
strings of Heisenberg field operators. This is perfectly justified for
particle accelerator experiments, where the background state is
essentially the vacuum. In early universe cosmology, however, the
background state is a thermal bath at temperature $T$, and a different
formalism is required.

\section{Finite-Temperature Green's Functions}

The finite-temperature Green's function for a scalar field $\varphi$ is
defined by replacing the vacuum expectation value with a thermal average
over the grand canonical ensemble,
\begin{equation}
	G_n^\beta(x_1,\ldots,x_n)
	= \frac{\mathrm{Tr}\left\{e^{-\beta H}\,
		T\left[\varphi(x_1)\cdots\varphi(x_n)\right]\right\}}
	{\mathrm{Tr}\,e^{-\beta H}},
\end{equation}
where $\beta = 1/T$ is the inverse temperature and the trace is taken
over all states of the system. This is the average value of the
time-ordered product of $n$ field operators in the grand canonical
ensemble with no chemical potential.

The finite-temperature generating functional is defined analogously,
\begin{equation}
	Z^\beta(J) = \frac{\mathrm{Tr}\left[e^{-\beta H}\,
			T\exp\!\left(i\int J(x)\,\varphi(x)\,d^4x\right)\right]}
	{\mathrm{Tr}\,e^{-\beta H}}.
\end{equation}
The essential difference between the zero-temperature and
finite-temperature formalisms lies in the boundary conditions on the
set of paths over which the functional integral is taken. At zero
temperature the transition amplitude is
\begin{equation}
	{}_H\langle\psi_1,t_1|\psi_2,t_2\rangle_H
	= \int\left[d\varphi(s,\mathbf{x})\right]
	\psi_1^*[\varphi(t_2,\mathbf{x})]\,\psi_2[\varphi(t_1,\mathbf{x})]\,
	\exp\!\left[i\int_{t_i}^{t_f}ds\int d^3x\,
		\mathcal{L}(\varphi,\partial\varphi)\right],
\end{equation}
with $t_i\to-\infty$ and $t_f\to+\infty$. At finite temperature the
boundary wavefunctionals $\psi^*[\varphi(t_f,\mathbf{x}_f)]\,
	\psi[\varphi(t_i,\mathbf{x}_i)]$ are replaced by
\begin{equation}
	\sum e^{-\beta E(\psi)}\,\psi^*[\varphi(t_f,\mathbf{x}_f)]\,
	\psi[\varphi(t_i,\mathbf{x}_i)],
\end{equation}
which implements the thermal weighting of states. It is also worth noting
that the finite-temperature effective potential is not gauge invariant:
only in physical gauges does the trace correspond to a sum over all
physical states of the system.

\section{Periodicity and the KMS Condition}

A fundamental consequence of the thermal boundary conditions is that the
finite-temperature Green's functions satisfy periodicity conditions in
Euclidean time. Working with Euclidean times $t\in[0,-i\beta]$, the
time-ordering is
\begin{equation}
	T\left[\varphi(x)\,\varphi(y)\right] =
	\begin{cases}
		\varphi(x)\,\varphi(y)         & \text{if } ix^0 > iy^0, \\
		(-1)^s\,\varphi(y)\,\varphi(x) & \text{if } ix^0 < iy^0,
	\end{cases}
\end{equation}
where $s=0$ for bosonic fields and $s=1$ for fermionic fields. To derive
the periodicity condition, we evaluate the two-point function at $x^0=0$
with $iy^0 > ix^0$,
\begin{align}
	\left(\mathrm{Tr}\,e^{-\beta H}\right)
	G_2^\beta(x-y)\big|_{x^0=0}
	 & = (-1)^s\,\mathrm{Tr}\left[e^{-\beta H}\,
		                        \varphi(y)\,\varphi(x)\right] \notag                               \\
	 & = (-1)^s\,\mathrm{Tr}\left[e^{-\beta H}\,
		                        \varphi(y^0;\mathbf{y})\,\varphi(0,\mathbf{x})\right] \notag       \\
	 & = (-1)^s\,\mathrm{Tr}\left[\varphi(0,\mathbf{x})\,
		                        e^{-\beta H}\,\varphi(y^0;\mathbf{y})\right] \notag                \\
	 & = (-1)^s\,\mathrm{Tr}\left[e^{-\beta H}\,
		                        \underbrace{e^{\beta H}\,\varphi(0,\mathbf{x})\,
			                        e^{-\beta H}}_{\varphi(-i\beta,\,\mathbf{x})}\,
		                        \varphi(y^0;\mathbf{y})\right] \notag                    \\
	 & = (-1)^s\,\mathrm{Tr}\left[e^{-\beta H}\,
		                        \varphi(-i\beta,\mathbf{x})\,\varphi(y^0;\mathbf{y})\right] \notag \\
	 & = (-1)^s\left(\mathrm{Tr}\,e^{-\beta H}\right)
	G_2^\beta(x-y)\big|_{x^0=-i\beta},
\end{align}
where in the third line we used the cyclicity of the trace, and in the
fourth line we used the Heisenberg equation of motion to identify
$e^{\beta H}\varphi(0,\mathbf{x})e^{-\beta H} = \varphi(-i\beta,\mathbf{x})$.
Dividing through by $\mathrm{Tr}\,e^{-\beta H}$,
\begin{equation}
	\boxed{G_2^\beta(x-y)\big|_{x^0=0}
		= (-1)^s\,G_2^\beta(x-y)\big|_{x^0=-i\beta}.}
\end{equation}
This is the Kubo-Martin-Schwinger (KMS) condition: bosonic Green's
functions are periodic in Euclidean time with period $\beta$, while
fermionic Green's functions are antiperiodic.

\section{Matsubara Frequencies and Feynman Rules}

The periodicity in Euclidean time has a direct consequence for the
momentum-space structure of the theory. Vertex factors are unaffected by
the change in boundary conditions, and the momentum-space propagator
remains of the same form. However, the periodicity forces the
energy component of the loop momentum to be discretized. The allowed
frequencies, known as Matsubara frequencies, are
\begin{align}
	\omega_n & = \frac{2n\pi}{\beta}
	\qquad\text{(Bose fields)},          \\
	\omega_n & = \frac{(2n+1)\pi}{\beta}
	\qquad\text{(Fermi fields)},
\end{align}
for integer $n$. The Feynman rules are modified accordingly. The
correspondence between zero-temperature and finite-temperature rules is
summarized in the following table.

\begin{center}
	\begin{tabular}{lcc}
		 & Zero temperature                                                     & Finite temperature \\[6pt]
		\hline                                                                                       \\
		Loop integral
		 & $\displaystyle\int\frac{d^4k}{(2\pi)^4}$
		 & $\displaystyle\frac{1}{-i\beta}\sum_{n=-\infty}^{\infty}
			   \int\frac{d^3k}{(2\pi)^3}$                                   \\[12pt]
		Vertex $\delta$-function
		 & $(2\pi)^4\,\delta^{(4)}\!\left(\sum k_i\right)$
		 & $\displaystyle\frac{\beta}{i}(2\pi)^3\,
			   \delta_{\sum\omega_n}\,\delta^{(3)}\!\left(\sum\mathbf{k}_i\right)$
		\\[6pt]
		\hline
	\end{tabular}
\end{center}

The replacement of the continuous energy integral by a discrete sum over
Matsubara frequencies is the defining feature of the imaginary-time
formalism. In the next section we use these rules to compute the
finite-temperature corrections to the effective potential.


